\section{Introduction}\label{sec:intro}

Quantum computations are intrinsically susceptible to noise and hence, quantum error correction (QEC) is essential for achieving fault tolerant quantum computing. The current landscape of the field is changing as it is transitioning from the noisy intermediate scale era towards more concrete demonstrations of fault tolerance. The current paradigm relies on qubit-based codes and architectures but an alternative approach has been proposed in the avenue of qudits. This paper will specifically look at the $q=3$ case of qudits, referred to as qutrits, and will work in the early fault tolerant quantum computing (early-FTQC) regime. At this stage, the goal is to realize small logical qutrits with protection against errors, before full scale QEC is feasible. 

Qutrits are three-level systems that present compelling alternatives for early FTQC research, where the aim is to study encoded logical devices before full error-correction cycles are practical. They display a larger and more connected computational state space than qubits: a single qutrit stores $\log_2 3 \approx 1.585$ bits of quantum information, increasing information density and, in principle, reducing both computational and ancillary overheads. Algorithmic primitives admit direct ternary advantages, for example, base-3 decompositions reduce \emph{(i)} circuit depth and \emph{(ii)} the count of costly entangling/non-Clifford gates in Shor’s modular exponentiation and generalized-Toffoli subroutines \cite{Gokhale2019}. Moreover, implementing the quantum Fourier transform (QFT) directly in \emph{base-3} reduces the register length from \(q_2=\lceil \log_2 N\rceil\) bits to \(q_3=\lceil \log_3 N\rceil\) trits; since a QFT over \(q\) digits uses \(\Theta(q^2)\) controlled rotations / controlled-\(\mathrm{SUM}\) operations, the entangling-gate count (and leading-order depth) drops by a factor \((\log_3 N/\log_2 N)^2 \approx 0.40\) at fixed \(N\), with explicit qudit-QFT circuit constructions given in \cite{PavlidisFloratos2021}. Similarly, Grover's search algorithm implemented with qutrits uses multi-level oracles and diffusion operators, reducing the gate count and circuit depth per iteration on hardware with native qutrit control \cite{Ivanov2012}. Beyond algorithms, higher-dimensional encodings improve quantum communication: in qutrit quantum key distribution (QKD), each carrier can encode more than one bit of key and protocols exhibit higher error tolerance than binary schemes \cite{BechmannPasquinucciPeres2000,Cerf2002}. Higher-dimensional systems strengthen foundational certificates used in quantum communication. With three-level systems, entangled pairs achieve larger violations of the Collins–Gisin–Linden–Massar–Popescu (CGLMP) Bell inequality than any qubit-based test, which increases the noise margin available to device-independent protocols \cite{Collins2002}. Further, qutrits are the smallest systems that can violate the Klyachko–Can–Binicioglu–Shumovsky (KCBS) contextuality inequality, improving robustness in semi-device-independent protocols \cite{Klyachko2008}. Qutrits also provide a natural language for simulating spin-1 and \(\mathrm{SU}(3)\) models and related phenomena, broadening the space of near-term scientific targets \cite{Senko2015}.
From a QEC standpoint, \emph{magic-state distillation} (MSD) in odd prime dimensions is especially attractive: quantum Reed Muller codes provide transversal qutrit gates enabling small block, high rate MSD protocols across all prime \(q\) \cite{CampbellPRX2012}, and working directly with \((q>2)\)-level systems further improves MSD overheads and provides routes to universal non-Clifford gates \cite{CampbellPRL2014}. Despite these advantages, systematic stabilizer-code design and logical benchmarking for qudits remain comparatively underexplored and merit systematic study. These considerations align with hardware reality: superconducting transmons, trapped ions, and photonic platforms natively provide multilevel structure.  Multi-platform demonstrations of coherent qutrit control, randomized benchmarking, and entangling operations now exist \cite{Morvan2021,Ringbauer2022,Goss2024}.

In this work, we employ the $[[5,1,2]]_3$ quantum error detection code to encode logical information. The data-qutrit layout and the logical gates are tailored to the connectivity of IQM’s superconducting Garnet QPU.
% TODO(authors): confirm which ancilla layout the Garnet coupler map supports (q1 at the center with one ancilla per stabilizer and a long-range q2-q3 coupler for the logical gates, or q2 at the center with q2 and q3 diagonal neighbors and two sites each for G1 and G4). The sentence above ("The data-qutrit layout and the logical gates are tailored to the connectivity of IQM's superconducting Garnet QPU") holds for the q2-centered patch; in the q1-centered one-ancilla layout, which leads the results, SWAP_{2,3} and CZ_{2->3} need a long-range q2-q3 coupler. Confirm or scope that sentence.
We consider two layouts of the code on a square grid, with ancillas at the four conventional surface-code sites between the data qutrits and check rounds that couple each ancilla only to its neighboring data qutrits (\autoref{sub:ancilla_layouts}). In the one-ancilla layout $q_1$ sits at the center, every stabilizer has one site next to its three data qutrits, and the logical gates need a long-range coupling between $q_2$ and $q_3$, which sit at opposite corners. In the two-ancilla layout $q_2$ sits at the center and $q_2$ and $q_3$ are diagonal neighbors, but no site is next to all three data qutrits of either boundary stabilizer, $G_1$ or $G_4$. This layout measures each of them with two sites that pass a partial syndrome through a shared data qutrit. We characterize the code and its constant-depth transversal logical Clifford gates with the \texttt{sdim} qudit stabilizer simulator \cite{sdim}, establishing a numerical baseline prior to hardware execution. Since the distance $d=2$ code supports error detection but not faithful correction, we adapt standard protocols to Logical Randomized Benchmarking with Detection (LRB-D), which isolates logical performance by postselecting on trivial syndromes. We simulate LRB-D under two circuit-level noise models of Gidney et al.~\cite{GidneyHoneycomb2021} that we generalize to qutrits: the superconducting-inspired SI1000 model and the uniform SD6 model, both parameterized by the two-qutrit gate error probability $p$. Every circuit is sampled through the exact detector error model of \texttt{sdim}, which reaches the $3\times 10^{7}$ shots per depth point that heavy postselection demands. Our objective is to compare the average fidelity of postselected logical Cliffords with that of physical single-qutrit Cliffords under the same noise model at the same $p$. We locate the pseudo-threshold $p^\star(k)$ below which the logical error is the smaller one, as a function of the interval $k$ between the check rounds whose syndromes enter the postselection. We do so for both layouts, each with nine qutrits.

% p* and their standard errors (seq_boot_se, sequence bootstrap) from LRB-experiment-data-slurm/comparisons/pseudo-thresholds-current.csv (one-ancilla rows) and pseudo-thresholds-conv-sites-current.csv (two-ancilla rows); layout gaps 1 - p*(two)/p*(one), each layout against the RB of its own run, with the relative standard errors of the two runs added in quadrature; SD6-minus-SI1000 gap differences: paper-rework/scratch/writer-conv/analysis.py (section S1)
With one ancilla per stabilizer and postselection after every forward logical Clifford ($k=1$), we find $p^\star(1) = 0.0292(6)$ under SI1000 and $p^\star(1) = 0.0399(15)$ under SD6, where the parentheses give one standard error from a bootstrap over the random Clifford sequences (\autoref{sub:pseudo_thresholds}).
Near these values almost every shot is rejected at larger depth. The pseudo-thresholds of both models therefore use additional runs at high $p$, at $p \geq 0.015$ under SI1000 and $p \geq 0.025$ under SD6, with $2\times 10^{7}$ shots per sequence and depth on the depths $0, 1, 2, 3, 4$ and $6$, of which the fits for $k \ge 2$ use the even ones (\autoref{sub:results_setup}).
At $p = 10^{-3}$ the logical error per unit depth is more than 90 times smaller than the physical one under both models. The SI1000 baseline is the harder one to beat, because there the one-qutrit gates of the physical baseline fail with probability $p/10$ rather than $p$. The pseudo-threshold decreases as fewer rounds enter the postselection, to $p^\star(4) = 0.0132(2)$ under SI1000 and $0.0155(4)$ under SD6.
These gains hold conditional on acceptance, and the acceptance probability falls quickly with depth. Under SI1000 at $p = 5\times 10^{-3}$ and $k = 1$, $63\%$ of shots are rejected at depth $m = 2$ and $98\%$ at $m = 10$.
The two-ancilla layout gives lower pseudo-thresholds at every $k$, by 9 to 18\% under SI1000 and by 21 to 33\% under SD6, where each layout is compared with the RB baseline of its own run and the gaps amount to 3.9 to 15.8 standard errors (\autoref{sub:layout_comparison}). The simulations charge the long-range $q_2$--$q_3$ gate of the one-ancilla layout like any other two-qutrit gate, so these gaps measure the cost of the two-ancilla check round and not that of the long-range coupling. The two-ancilla check round has 18 two-qutrit gates in 11 gate moments, against 12 gates in 7 gate moments with one ancilla, which raises the expected number of faults per round from $67.9p$ to $76.3p$ under SI1000 and from $73p$ to $103p$ under SD6, where idle qutrits fail with probability $p$. The gap is larger under SD6 than under SI1000 at $k = 1$, $3$ and $4$, by 3.6 to 5.8 standard errors, while at $k = 2$ the difference amounts to 2.7 standard errors.

% Qubit comparison: qubit p* from comparisons/pseudo-thresholds-qubit-current.csv (qubit-one-ancilla rows) and pseudo-thresholds-qubit-conv-sites-current.csv; ratios with relative standard errors in quadrature, RB and LRB-D error ratios at equal p, word lengths: paper-rework/scratch/writer-qubit/analysis.py (sections R1, B1, L1); LRB-D ratio window and RB estimates rechecked in paper-rework/scratch/corrector-qubit-1/corrector.py (sections A, B)
We also repeat the production simulations for the qubit $[[5,1,2]]$ code, with the same check rounds, logical gates and fits and the qubit forms of the same noise models, and replace the high-$p$ extension by near-threshold refinement runs (\autoref{sub:qubit_comparison}). Its pseudo-thresholds are 0.29 to 0.38 times the qutrit values under SI1000 and 0.45 to 0.56 times under SD6, at $k = 1$ to $4$ and in both layouts. Two effects combine. A qubit Clifford averages 2.54 primitive gates against 5.17 for a qutrit, and at the same $p$ a noisy qubit gate has $8/9$ of the average error of a noisy qutrit gate. The physical qubit RB error is 0.42 to 0.49 times the qutrit one, close to the first-order estimate $(8/9)(2.54/5.17) = 0.44$ from these two factors, and higher orders account for the larger values under SD6 at large $p$. The postselected logical error per unit depth is instead larger for the qubit: at the same $p$ between $10^{-3}$ and $10^{-2}$ it is 2.4 to 3.0 times the qutrit one under SI1000 and 1.3 to 2.1 times under SD6, and we have not identified the mechanism behind it.

Logical dynamics can be non-Markovian even under Markovian physical noise, because the encoded state couples to syndrome subspaces in which undetected errors persist and act on later logical gates \cite{Ziyad2025}. Intermediate postselection discards the trajectories in which a check round detects such an excursion. Our simulations are consistent with this picture, but they do not isolate the mechanism. With terminal postselection alone, the accepted logical signal is not a single exponential in depth: it decays slowly at small depth and then steepens. At depth $m = 22$ it drops below the physical RB signal for $p \geq 10^{-2}$ under SI1000 and for $p \geq 5\times 10^{-3}$ under SD6. Under SI1000 a single exponential describes the accepted decays for $k = 1$ to $4$, closely for $k \le 2$ and approximately for $k = 3$ and $4$, where the traces show small steps at depths that are not multiples of $k$ (\autoref{sub:interval_results}). Under SD6 the accepted decay already steepens with depth at $k = 1$, and under both models the $\chi^2$ of the fit tends to grow with $k$ (\autoref{sub:interval_results}). For $p \geq 0.02$ and $k \le 4$ the exponential form can be tested only over short depths, because almost all shots at large depth are rejected. Each fit of the production runs there rests on two to five depths, and the fits of the high-$p$ runs use at most six depths at $k = 1$ and at most the four even depths at $k \ge 2$, all with $m \le 6$.